% Drawing 74/05-1 -- folder 74, batch 05, archive page 81 (PDF page 82).
% Transcription: Dans la marge gauche, des schémas de points (gros points, petits points, croix) et de segments pleins ou pointillés : un segment et une demi-diagonale pointillée en haut, puis deux cadres de $3\times 3$ points reliés par des flèches, illustrant la position d'un carré et de son centre dans le pavé.
% Kind: marginal point schemas (squares and their centres). Method 2 (schematic redraw from the facsimile). Status: partial.
% Notes: Left margin, top to bottom. A bold segment, a bold stroke from its right end down to the left, and a dotted half-diagonal under it. A square with its centre marked by a point; the letters y, x' (top) and x, y' (bottom) at its corners (read on the facsimile, not in the transcription). Far below, two rounded frames of points, each with two columns of three points (the middle row big dots, the others small; a small cross by the top right one of the lower frame); between them, two big points and a small one under the upper frame, a group of six points (big dots, a cross) to the right, and arrows: one down the left side from the upper frame to the lower, one from the upper frame towards the group, one from the group back down. Every point is placed as on the page. Struck out and omitted: the scribble of short strokes under the upper frame. Still differs: the small faint marks around the dotted half-diagonal are not drawn; where the arrows end is approximate. Redrawn 2026-11-09 (issue 128): to check again.
% Coordinates: 1 unit = 1 mm = 1/100 of the crop width (bbox in page fractions: [0.08, 0.27, 0.24, 0.88]).
% Redrawn by Claude (claude-opus-5-5), issue #73, 2026-11-09; not checked by a human; not compiled (no TeX in the session).
% Rights: reproduces part of the work -- behind RENDER_TRANSCRIPTS, never published by itself.
\documentclass[tikz,border=3pt]{standalone}
\usepackage{amsmath,amssymb}
\usetikzlibrary{patterns,arrows.meta}
\tikzset{pencil/.style={black!45}, struck/.style={black!35}, mark/.style={circle,draw,inner sep=1.3pt}, pt/.style={circle,fill,inner sep=0.9pt}, opt/.style={circle,draw,fill=white,inner sep=1.1pt}}
\begin{document}
\begin{tikzpicture}[x=1mm,y=1mm,line cap=round,line join=round,font=\scriptsize]
\draw[very thick] (37,-18) -- (60,-17);
\draw[very thick] (60,-17) -- (35,-32.5);
\draw[densely dotted] (58,-22) -- (40,-42);
\draw (31,-164) -- (58,-163);
\draw (58,-163) -- (59,-195);
\draw (59,-195) -- (32,-195);
\draw (32,-195) -- (31,-164);
\node[pt] at (46,-179) {};
\node[inner sep=1pt] at (25,-160) {\small $y$};
\node[inner sep=1pt] at (66,-159) {\small $x'$};
\node[inner sep=1pt] at (25,-202) {\small $x$};
\node[inner sep=1pt] at (69,-200) {\small $y'$};
\draw (22,-414) -- (40,-414) -- (43,-418) -- (43,-443) -- (40,-446) -- (22,-446) -- (19,-443) -- (19,-418) -- cycle;
\node[pt] at (24,-419) {};
\node[pt] at (35,-419) {};
\node[pt,very thick] at (24,-427) {};
\node[pt,very thick] at (36,-427) {};
\node[pt] at (25,-436) {};
\node[pt] at (36,-436) {};
\node[pt,very thick] at (20,-450) {};
\node[pt,very thick] at (30,-450) {};
\node[pt] at (39,-448) {};
\node[pt] at (56,-468) {};
\node[pt] at (71,-467) {};
\node[pt] at (83,-468) {};
\draw (53.7,-478.7) -- (56.3,-481.3);
\draw (53.7,-481.3) -- (56.3,-478.7);
\node[pt] at (69,-479) {};
\node[pt] at (83,-480) {};
\draw[-{Stealth[length=2mm]}] (19,-455) -- (18.5,-498);
\draw[-{Stealth[length=2mm]}] (39,-453) -- (43,-465);
\draw[-{Stealth[length=2mm]}] (42,-480) -- (37,-494);
\draw (20,-502) -- (41,-502) -- (44,-505) -- (44,-525) -- (41,-528) -- (20,-528) -- (17,-525) -- (17,-505) -- cycle;
\node[pt] at (22,-508) {};
\draw (33.8,-505.8) -- (36.2,-508.2);
\draw (33.8,-508.2) -- (36.2,-505.8);
\node[pt,very thick] at (22,-516) {};
\node[pt,very thick] at (34,-516) {};
\node[pt] at (23,-523) {};
\node[pt] at (34,-523) {};
\end{tikzpicture}
\end{document}
